\section*{Chapter \ref{ch:g10:project-web-application} --- Project: A Small Web Application}

\begin{solution}{exo:g10:project-web-application:1}
\code{/}: \code{home(request)}. \code{/book/12} and \code{/book/12/}: \code{book(request, "12")}, since the slashes at
both ends are stripped before splitting. \code{/book/12/loans}: three parts, one too many, so
\code{not_found(request)}. \code{/books/12}: no \omterm{def:g10:project-web-application:route}{route} \code{books}, so \code{not_found(request)}.
\end{solution}

\begin{solution}{exo:g10:project-web-application:2}
The title is \textit{Not found} and the paragraph is

\smallskip\noindent\hspace*{2em}\code{<p class="error">No page at /a&amp;b&lt;i&gt;</p>}

\smallskip\noindent The browser displays \textit{No page at /a\&b<i>}, as
typed, in the colour of the class \code{error}; no italic starts, since \code{<i>} was escaped.
\end{solution}

\begin{solution}{exo:g10:project-web-application:3}
A function \code{about(request)} that returns \code{app_m2.page("About", "<p>...</p>")}, and an entry
\code{"about": (about, 0)} in the \omterm{def:g10:project-web-application:route}{route} table (\code{exercises/app_ex.py}).
\end{solution}

\begin{solution}{exo:g10:project-web-application:4}
One for the home page, one for the book, four for the missing pages, one for the style sheet, one for the search,
two for the reservation (the \omterm{def:g10:http-clients-servers:methods}{POST}, then the \omterm{def:g10:http-clients-servers:methods}{GET} that follows the 303) and two refused reservations: $12$ requests.
\end{solution}

\begin{solution}{exo:g10:project-web-application:5}
\code{book} returns \code{None}, and so does \code{dispatch}. In the server, the line
\code{status, kind, text, headers = app.dispatch(request)} raises a \code{TypeError}: the server prints the
traceback, closes the connection without any answer, and goes on serving the other requests. The client sees the
connection closed without a response (\code{RemoteDisconnected} with \code{http.client}).
\end{solution}

\begin{solution}{exo:g10:project-web-application:6}
\code{author(request, name)} keeps the books whose \code{author} is \code{name}, and returns
\code{not_found(request)} if there are none; the \omterm{def:g10:project-web-application:route}{route} is \code{"author": (author, 1)}. The space in the name
arrives as \texttt{\%20} and the server decodes it. Zoe Stoneton has $7$ books, the most of any author
(\code{exercises/app_ex.py}).
\end{solution}

\begin{solution}{exo:g10:project-web-application:7}
The form's fields travel in the body of the \omterm{def:g10:http-clients-servers:methods}{POST}, not in the \omterm{def:g10:http-clients-servers:url}{query string}, so \code{request["query"]} is empty: the
book number is \code{""}, and every reservation gets \code{400} with the page \textit{No such book}. The test
\code{test_reserve_then_redirect} fails, since it expects the book's page after the redirection.
\end{solution}

\begin{solution}{exo:g10:project-web-application:8}
Read \code{key = request["query"].get("sort", "title")}, answer 400 if it is not one of the three, and list the books
in the order \code{sorted(BOOKS, key=lambda n: (BOOKS[n][key], n))}; the number \code{n} breaks ties, so the order
is the same at every request (\code{exercises/app_ex.py}).
\end{solution}

\begin{solution}{exo:g10:project-web-application:9}
The home page, the style sheet, one page per book and one \omterm{def:g10:data-tables:json}{JSON} answer per book: $2n + 2$. For $80$ books, $162$
pages (with \omterm{def:g10:http-clients-servers:methods}{GET}, \code{/reserve} answers 400, not 200). The addresses are far more numerous, without end: as in
exercise \ref{exo:g10:project-web-application:1}, a trailing slash or leading zeros (\code{/book/7/}, \code{/book/07},
\code{/book/007}, \dots) give another address for the same page.
\end{solution}

\begin{solution}{exo:g10:project-web-application:10}
The failing test raises an exception, the last line is never reached, and the server keeps running after the tests.
The next run cannot open the same port, and its server stops at once, while the tests talk to the old server, which
runs the old code. The fix is to stop the server in a \code{finally} block, or in \code{tearDownClass} as in
milestone~4, which \code{unittest} calls whatever the tests did.
\end{solution}

\begin{solution}{exo:g10:project-web-application:11}
A dictionary \texttt{METHODS = \{"reserve": \{"POST"\}\}}, with \texttt{\{"GET"\}} for every other \omterm{def:g10:project-web-application:route}{route}; before calling
the \omterm{def:g10:project-web-application:route}{route}'s function, \code{dispatch} compares \code{request["method"]} with the allowed set and otherwise answers
405 with the header \code{Allow}, which a 405 answer must carry, listing the allowed methods
(\code{exercises/app_ex.py}).
% ledger: g10:http-405-allow
\end{solution}

\begin{solution}{exo:g10:project-web-application:12}
With $80$ books and $20$ per page there are $\lceil 80/20 \rceil = 4$ pages. The page $p$ shows the books at
positions $20(p-1)$ to $20p - 1$ of the sorted list, the slice \code{order[(p - 1) * 20:p * 20]}. Pages $0$ and
$5$, and a page that is not a number, get a 404. The test requests pages $0$ to $5$, checks the statuses
$404, 200, 200, 200, 200, 404$, and counts $80$ links over the four pages (\code{tests/test_web_app.py}).
\end{solution}

\begin{solution}{exo:g10:project-web-application:13}
If the \omterm{def:g10:http-clients-servers:methods}{POST} is answered with a page, reloading repeats the \omterm{def:g10:http-clients-servers:methods}{POST}: $1 + 3 = 4$ reservations of the same book by the
same member. With a 303, the page shown is the answer to a \omterm{def:g10:http-clients-servers:methods}{GET}, and reloading repeats the \omterm{def:g10:http-clients-servers:methods}{GET}: $1$ reservation.
Browsers warn before repeating a \omterm{def:g10:http-clients-servers:methods}{POST} because a \omterm{def:g10:http-clients-servers:methods}{POST} may change data on the server, while a \omterm{def:g10:http-clients-servers:methods}{GET} must not.
\end{solution}

\begin{solution}{exo:g10:project-web-application:14}
Before recording, count the member's reservations over all the books, and answer 409 if there are already three. The
test sends four reservations for member $5$ and expects $303, 303, 303, 409$ (\code{exercises/app_ex.py}).
\end{solution}

\begin{solution}{exo:g10:project-web-application:15}
Thread A counts member $5$'s reservations: $2$, fewer than $3$. Thread B counts them too: still $2$. A records its
reservation ($3$), then B records its own ($4$). Each check was right when it was made, but the data changed between
the check and the recording. The fix is to make the check and the recording one step that no other thread can
interrupt: a lock (\code{threading.Lock}) that a thread holds from the count to the recording
(\code{exercises/app_ex.py}). If several server programs share the data, a lock inside one of them is not enough;
the database that holds the data must provide it, as a later chapter shows.
\end{solution}
